\documentclass[11pt,letterpaper]{article}

\usepackage[utf8]{inputenc}
\usepackage[T1]{fontenc}
\usepackage{amsmath,amssymb,amsthm}
\usepackage{graphicx}
\usepackage{booktabs}
\usepackage[margin=1in]{geometry}
\usepackage{hyperref}

\newenvironment{code}
  {\par\medskip\ttfamily}
  {\par\medskip}

\title{Paper Title Goes Here}
\author{Author Name \\ Affiliation}
\date{\today}

\begin{document}

\maketitle

\begin{abstract}
The $k$-color Ramsey number of a poset $P$ is the smallest dimension $d$ where for any $k$-coloring of the $d$ dimensional Boolean lattice, there is a monochromatic induced copy of $P$. We will describe this formalization in the Lean 4 proof assistant, where we define and rigorously prove one of these bounds for the 2-diamond.
\end{abstract}

\section{Introduction}
We will consider Ramsey bounds for posets. Denote $Q_n$ the $n$-dimensional Boolean lattice, the lattice containing all subsets of $[n]$ ordered by inclusion.
We define the $k$-color Ramsey number of a finite poset $P$ by the smallest dimension $d$ where in any $k$-coloring $\chi : \mathcal P([n]) \to [k]$ there exists a monochromatic induced copy of $P$. 


\section{Defining Basic Structures}

In this section we will briefly provide an overview of the main argument, where we present key definitions and lemmas. All code in this section was written by humans and are located in \texttt{Defs.lean}. Even though Lean's Mathlib includes partially ordered sets and directed graphs, not enough features are in these packages. In addition, when defining the boolean lattice $Q_n$ much of the proof relies on individual elements as sets, so custom definitions are provided.

We will first formalize the notion of a boolean lattice in Lean. A partial order is defined on $\mathcal P([n])$ to give us the notion of $Q_n$.
With this framework we are able to define $n$-diamonds and $n$-wedges. An order is placed on these wedges and basic facts about these orders are presented in the file.
\begin{code}
abbrev BooleanLattice (n : $\mathbb N$) : Type := Finset (Fin n)
    
def Diamond (n : $\mathbb N$) := Fin (n + 2) \\
instance (n : $\mathbb N$) : Fintype (Diamond n) := inferInstanceAs (Fintype (Fin (n + 2)))

def Wedge (n : $\mathbb N$) := Fin (n + 1) \\
instance (n : $\mathbb N$) : Fintype (Wedge n) := inferInstanceAs (Fintype (Fin (n + 1)))
\end{code}

\subsection{Defining Ramsey Colorings}

We will now construct the proposition where $P$ is an induced copy in $Q$ from first definitions. As $P$ and $Q$ are partial orders, we are able to construct an injective order-preserving function $f : P \to Q.$ Furthermore, we will define a $k-$coloring of a poset $P$ as simply a function $\chi : P \to [k],$ and are able to provide a proposition stating whether if given a coloring $\chi,$ there is a color $c$ where $p \in P$ implies $\chi(p) = c.$
\begin{code}
abbrev Coloring (k : $\mathbb N$) (V : Type) : Type := V $\to$ (Fin k)

def IsInducedCopy {P Q : Type} [PartialOrder P] [PartialOrder Q] (f : P → Q) := \\
  Function.Injective f $\wedge$ $\forall$ (a b : P), a $\leq$ b $\leftrightarrow$ f a $\leq$ f b


  def HasMonochromeSubposet (P Q : Type) [PartialOrder P] [PartialOrder Q] (k : $\mathbb N$) \\
    ($\chi$ : Coloring k Q) := \\
  $\exists$ (c : Fin k), $\exists$ (f : P $\to$ Q), IsInducedCopy f $\wedge$ $\forall$ (x : P), $\chi$ (f x) = c
  
\end{code}

It is now possible to construct definitions of upper and lower Ramsey bounds, enabling us to define a multicolor Ramsey bound. We say $R_k(P) \leq n$ if all $k$-colorings of $Q_n$ provide us a monochromatic induced $P$, and similarly $R_k(P) < m$ if we are able to find a coloring $\chi$ where there is no monochromatic induced $P$ in $Q_m.$ We are then able to translate these definitions in Lean. Two smaller definitions definining Ramsey upper and lower bounds are provided, and they are combined to form our main definition.


\section{Results}
Present your findings, theorems, or empirical measurements.

\section{Conclusion}
Summarize key takeaways and future work.

\bibliographystyle{plain}
\bibliography{references}

\end{document}